% =============================================================================
% Дүгнэлт (дугааргүй)
% =============================================================================

\defformat
\addchap{Дүгнэлт}
\label{ch:conclusion}
\pagecolor{white}

Энэхүү ажлаар нэг GPU-д ажиллах нээлттэй эхийн Arctic-Text2SQL-R1-7B загварт суурилсан, банк бус санхүүгийн
байгууллагын өгөгдлийн сан дээр монгол хэлээр асуусан асуултыг SQL болгож хариулдаг системийг бүтээж, бүрэлдэхүүн бүрийн
нөлөөг хяналттай туршилтаар хэмжив.

Шинэ асуултаас бүрдсэн test дээр нэмэлт сургалттай суурь (D2) 21.0\% байсан бол орчуулга ба засвартай H4 34.0\%
($p = 8.6\times10^{-4}$), 8 хувилбарын санал хураалттай E2 38.5\%, тодруулгатай E3c 46.5\% хүрэв. Мэдлэг хайлт (+12.0
пункт) ба орчуулга (+11.3) хамгийн их, өөрийн DB-ийн жишээгүй нэмэлт сургалт (+1.4) ба алдааны мессежээр засах (+0.3)
хамгийн бага нэмэр өгсөн. Алдааны 38\% нь орчуулгаас үүдсэн тул санхүүгийн хэллэгийн орчуулгын өгөгдөл бэлтгэж,
орчуулагчийг нэмж сургав: санхүүгийн 8 500 мөрөөр QLoRA сургасан MADLAD-400-3B-ийн chrF++ 58.0 → 74.8 болж, түүгээр
орчуулсан нэг үүсгэлттэй тохиргоо (M1a) dev-ийн хариулах боломжтой асуултад 56.7\% буюу Qwen орчуулгатай I3t-ээс 9.8 пункт
өндөр гарсан (test дээр баталгаажуулах шаардлагатай). Сэдэв батлах даалгаврын зорилтуудын биелэлтийг Хүснэгт~\ref{tab:goals-met}-д
нэгтгэв.

\begin{table}[H]
\centering\small
\caption{Зорилгын биелэлтийн нэгтгэл}
\label{tab:goals-met}
\begin{tabularx}{\textwidth}{p{0.5cm}X p{5.4cm}}
\toprule
№ & Зорилт & Биелэлт \\
\midrule
1 & Мэдлэгийн сан (RAG) бүрдүүлэх & 28 view, 425 багана, 45 нэр томьёо, 20 KPI, 635 жишээ (§\ref{sec:db}) \\
2 & Prompt ба pipeline боловсруулах & 7 алхамтай дамжлага, 84 нэрлэсэн тохиргоо (Бүлэг~\ref{ch:method}, \ref{ch:system}) \\
3 & Монгол хэл ба өөрийн DB-д тохируулан сургах & C1--C3, D2 (A100), D3 (T4); орчуулагчийн M1 сургалт (chrF++ 58.0 → 74.8); H3 батлагдаагүй \\
4 & Олон хувилбараас сонгох, засах, татгалзах & E1, E2, E3c, E4; test хариулах боломжтой EX +17/−3 ($p = 0.0026$) \\
5 & Тестийн багцаар хэмжиж тайлагнах & dev 300 + test 200; test 21.0\% → 46.5\% (§\ref{sec:numbers}) \\
\bottomrule
\end{tabularx}
\end{table}

\textbf{Оруулсан хувь нэмэр.} (1) Монгол хэлний Text2SQL-ийн 500 асуулттай үнэлгээний сан; (2) бага нөөцтэй хэл ба 7B
загварт мэдлэг хайлт ба орчуулга хамгийн их нөлөөтэйг McNemar тестээр харуулсан эмпирик нотолгоо; (3) алдааны эх
үүсвэрийн хэмжилт ба орчуулгын өгөгдөл; (4) үнэгүй үүлэн квот дээр давтагдахуйц систем ба Text2SQL Studio хэрэгсэл.

\textbf{Цаашдын чиглэл.}
\begin{enumerate}
  \item Сургасан орчуулагчтай M1a тохиргоог test дээр нэг удаа хэмжиж, SQL бичигчийн (D3, схемд тогтвортой INV1) сургалтын
        үр дүнг dev дээр шийдэх; M1-ийн DB утга хадгалах сул талыг (13.6\%) засах;
  \item Үнэлгээний санг хоёр домэйн шинжээчээр шалгаж, бүх харьцуулалтыг нэг төхөөрөмж дээр давтах;
  \item Латин асуулт дээр тусад нь хэмжиж, бодит өгөгдлийн хамгаалалттай хуулбар дээр shadow горимоор туршиж хэрэглэгчийн санал авах.
\end{enumerate}

\textbf{Нийгмийн нөлөө.} Монгол хэлээр өгөгдлийн сангаас тоо мэдээ авах боломж нь SQL мэдэхгүй ажилтнуудын шинжээчээс
хамаарах хамаарлыг бууруулж, шийдвэр гаргалтыг хурдасгана. Нээлттэй эхийн загвар ба нэг GPU-д суурилсан тул жижиг, дунд
байгууллага ч өгөгдлөө гадагш гаргахгүйгээр ашиглах боломжтой.
